\section{Scope, Artifacts and Conclusion}
The neural results describe one graph, an assumed RF family, a synthetic plant, unmatched healthy reserve and sparse failures. Walk-only contact by 0.8\,s follows from this plant's speed construction, not connectome-specific advantage. Functional row deletions establish no biological tolerance, physical upset rate, hardware safety, energy or area benefit. The interface makes declared pairing, timing and exposure testable across two native schemas; it does not prove arbitrary fault coverage or authentic neural execution.

The fixed V10 code/data snapshot\cite{artifactV10} provides both adapters, 576 neural test records, 96 tank examples, frozen configurations and a minimal replay entry point. Original generation-time versions are preserved. Saved-data replay requires no external matrices; fresh neural simulation requires separately obtained, hash-matched inputs. Third-party tables with unresolved redistribution terms are excluded. A fresh checkout and clean dependency installation previously reproduced 448 V9 saved traces on the same host, without establishing second-host neural replication. New RF tests used four Xeon W-2223 workers, about 179\,s median episode wall time and at most 1.21\,GiB lifetime working set per worker.

Paired connection loss produced different collision outcomes despite equal healthy success. Smaller, uncertain contrasts in separately recalibrated RF realizations and the observed warning errors qualify the baseline finding. Executable pairing checks and independent native replay make these conditional results auditable, with limited reuse beyond the fly task. This task-level assessment provides a basis for evaluating how future device-level nonidealities propagate through neuromorphic controllers to closed-loop behavior.
